\documentclass[10pt,letterpaper]{article}

\usepackage[utf8]{inputenc}
\usepackage[T1]{fontenc}
\usepackage[spanish,es-noshorthands]{babel}
\usepackage{lmodern}
\usepackage[top=1.5cm,bottom=1.5cm,left=1.8cm,right=1.8cm]{geometry}
\usepackage{xcolor}
\usepackage{enumitem}
\usepackage[hidelinks]{hyperref}

\definecolor{notiuCrimson}{HTML}{AA0330}
\definecolor{notiuDark}{HTML}{260A13}
\definecolor{notiuWarmBg}{HTML}{F9F5F3}

\pagestyle{empty}
\setlength{\parindent}{0pt}
\setlength{\parskip}{3.5pt}

\begin{document}

% Encabezado compacto
\noindent
\colorbox{notiuDark}{%
  \parbox{\dimexpr\textwidth-2\fboxsep\relax}{%
    \vspace{4pt}
    {\Large\bfseries\textcolor{white}{NOTIU --- Términos y Condiciones de Uso (Versión Preliminar)}}\\[2pt]
    {\footnotesize\textcolor{white}{Versión 0.5.0 \textbar\ Vigencia: Noviembre de 2026 \textbar\ Proyecto Académico --- Construcción de Software}}
    \vspace{4pt}
  }%
}

\vspace{4pt}
{\small
El presente documento establece las condiciones básicas de uso de \textbf{NotiU}, aplicación multiplataforma de organización académica y gestión de notificaciones estudiantiles. Al registrarse o utilizar el servicio en cualquiera de sus entornos de prueba, el usuario acepta los siguientes términos:}

\vspace{2pt}
\begin{enumerate}[leftmargin=1.4em, itemsep=3.5pt, topsep=2pt, parsep=0pt, label=\textbf{\textcolor{notiuCrimson}{\arabic*.}}]

  \item \textbf{Naturaleza Académica y Objeto del Servicio.}
  NotiU es un proyecto de software en fase de prototipo desarrollado con fines estrictamente formativos en la asignatura \emph{Construcción de Software}. Su propósito es facilitar a los estudiantes universitarios la consulta centralizada de avisos académicos, el seguimiento de materias, horarios, apuntes, recordatorios y la proyección orientativa de calificaciones.

  \item \textbf{Registro de Cuenta y Uso Aceptable.}
  Para acceder a las funcionalidades personalizadas, el usuario debe registrarse proporcionando información veraz (nombre, correo electrónico y contraseña de mínimo 8 caracteres). El usuario es responsable de resguardar sus credenciales y tokens de sesión, comprometiéndose a emplear la plataforma únicamente con fines personales, educativos y lícitos, absteniéndose de vulnerar cuentas de terceros, alterar la API o automatizar peticiones abusivas.

  \item \textbf{Fase de Prototipo y Almacenamiento Temporal.}
  En su versión actual (\texttt{v0.5.0}), el backend opera como entorno demostrativo con persistencia temporal en memoria (\emph{in-memory}). En consecuencia, las cuentas, materias, notas y configuraciones registradas pueden restablecerse cuando el servidor se reinicie o actualice. El servicio se ofrece ``tal cual'' (\emph{as is}), sin garantías de disponibilidad continua ni respaldo permanente de información durante esta etapa experimental.

  \item \textbf{Privacidad y Tratamiento Básico de Datos.}
  NotiU recolecta únicamente los datos necesarios para su funcionamiento académico: identificación básica, perfil universitario (institución, carrera, semestre, intereses), asignaturas, calificaciones y preferencias de alertas o visualización. Las contraseñas se almacenan protegidas mediante funciones criptográficas de derivación (\texttt{scrypt}) y nunca en texto plano. Siguiendo los principios de la Ley 1581 de 2012 (Colombia), los datos no se comercializan ni se ceden a terceros; el usuario puede restringir la visibilidad de su perfil ante otros estudiantes desde los ajustes de privacidad o solicitar la eliminación de sus datos.

  \item \textbf{Alcance Informativo e Integraciones Externas.}
  Los cálculos de promedios ponderados por cortes, estimaciones de nota mínima necesaria para aprobar, temporizadores de estudio y notificaciones vinculadas con plataformas externas (tales como Moodle, Google Classroom, Microsoft Teams o correo electrónico) constituyen herramientas de apoyo personal. Dichos resultados tienen carácter meramente orientativo y no sustituyen las calificaciones, fechas límite ni certificaciones oficiales emitidas por cada institución educativa.

  \item \textbf{Propiedad Intelectual.}
  El código fuente, la arquitectura del sistema, la interfaz de usuario, la denominación \textbf{NotiU} y el isotipo del zorro con punta de estilógrafo son titularidad de su equipo desarrollador (\emph{Alejandro Rodríguez, Alejandro Martínez, David Muñoz, Yorman Sanabria y Jireth Suárez}), sin perjuicio de las bibliotecas y herramientas de código abierto empleadas bajo sus respectivas licencias.

  \item \textbf{Modificaciones y Contacto.}
  Estos términos preliminares podrán ajustarse conforme el proyecto incorpore base de datos persistente, autenticación con proveedores externos y despliegue móvil en producción. Cualquier consulta, reporte técnico o solicitud relacionada con el proyecto puede canalizarse a través del repositorio oficial: \url{https://github.com/Manum03/NotiU}.

\end{enumerate}

\vfill
\noindent\textcolor{notiuCrimson}{\rule{\textwidth}{0.6pt}}\\[2pt]
{\scriptsize\textcolor{gray}{NotiU (\texttt{v0.5.0}) --- Curso de Construcción de Software, Docente: Diego Andrés Baquero Tibocha. Documento preliminar de una (1) página para validación funcional y académica.}}

\end{document}
